\section{Related Works}

Learning a shared hypothesis space from related tasks is a classical
formulation of inductive-bias learning~\cite{baxter2000model}.
For neural networks, structural constraints can express locality,
equivariance and parameter sharing~\cite{ravanbakhsh2017equivariance}.
Deep Sets~\cite{zaheer2017deep} provides a useful example of explicitly
encoded invariance, whereas here the relevant connections and equality
groups are not supplied to the learner.

The closest structural competitors parameterize weights as a shared
linear map applied to task-specific coefficients. Yeh
et al.~\cite{yeh2022equivariance} learn an assignment matrix using a
differentiable relaxation and an implicit hypergradient. MSR
\cite{zhou2021meta} meta-learns a continuous reparameterization,
together with numerical initialization and adaptation rates.
We compare native continuous representations separately from hard
indexators, and distinguish transfer of a structural constraint from
transfer of learned numerical initial values.

Weight tying and quantization also learn equality groups.
Soft weight tying~\cite{nowlan1991adaptive}, trained quantization
\cite{han2016deep}, LUT-Q~\cite{cardinaux2018iteratively} and
weight-sharing regularization~\cite{shakerinava2024weight} establish
that learning a dictionary and assignments is not itself new.
Our experimental question is whether a common assignment extracted
from several source solutions remains useful after replacing every
task-specific numerical value. The proposed component is joint
calibrated projection of aligned solutions, followed by functional
comparison of a finite bank; the dictionary representation and the
general principle of learning structure across tasks are prior art.
